\documentclass[11pt,a4paper]{article}

\usepackage[margin=0.85in,headheight=15pt]{geometry}
\usepackage{amsmath,amssymb,amsthm}
\usepackage{fancyhdr}
\usepackage{tcolorbox}
\usepackage{microtype}
\usepackage{lastpage}
\usepackage[colorlinks=true,linkcolor=blue,urlcolor=blue]{hyperref}

% Header / Footer setup
\pagestyle{fancy}
\fancyhf{}
\lhead{\textbf{Aryan Maurya} \quad (\textbf{Roll No:} 24220MAT051)}
\rhead{\textbf{Abstract Algebra Assignment}}
\cfoot{Page \thepage\ of \pageref{LastPage}}
\renewcommand{\headrulewidth}{0.4pt}

\tcbset{
    colback=gray!6!white,
    colframe=black!75!white,
    fonttitle=\bfseries,
    coltitle=white,
    rounded corners,
    boxrule=0.8pt,
    arc=2mm,
    top=6pt,
    bottom=6pt,
    left=10pt,
    right=10pt
}

\begin{document}

\begin{center}
    {\LARGE \textbf{Department of Mathematics}}\\[3pt]
    {\large \textbf{Abstract Algebra / Group Theory}}\\[6pt]
    \rule{\textwidth}{0.8pt}
    \vspace{2pt}
    \begin{tabular*}{\textwidth}{l @{\extracolsep{\fill}} r}
        \textbf{Student Name:} Aryan Maurya & \textbf{Roll Number:} 24220MAT051 \\
        \textbf{Topic:} Center of Dihedral Groups $Z(D_{2n})$ & \textbf{Date:} \today
    \end{tabular*}
    \vspace{-4pt}
    \rule{\textwidth}{0.8pt}
\end{center}

\vspace{4pt}

\begin{tcolorbox}[title={Problem Statement (Chapter 2, Problem 7)}]
Let $n \in \mathbb{Z}$ with $n \ge 3$. Prove the following:
\begin{enumerate}
    \item[\textbf{(a)}] $Z(D_{2n}) = 1$ if $n$ is odd.
    \item[\textbf{(b)}] $Z(D_{2n}) = \{1, r^k\}$ if $n = 2k$.
\end{enumerate}
\end{tcolorbox}

\vspace{4pt}

\section*{Complete and Rigorous Proof}

\subsection*{1. Preliminaries and Presentation of $D_{2n}$}
The dihedral group $D_{2n}$ of order $2n$ is the symmetry group of a regular $n$-gon. It is generated by a rotation $r$ of order $n$ and a reflection $s$ of order $2$, with the standard presentation:
\[
D_{2n} = \big\langle r, s \;\big|\; r^n = 1,\; s^2 = 1,\; sr = r^{-1}s \big\rangle.
\]

From the relation $sr = r^{-1}s$, an easy induction on $i \ge 0$ yields the general commutation formula:
\begin{equation} \label{eq:rel}
s r^i = r^{-i} s \quad \text{for all } i \in \mathbb{Z}.
\end{equation}

Every element $z \in D_{2n}$ can be uniquely expressed in canonical form as:
\[
z = r^i s^j \quad \text{where } 0 \le i < n \text{ and } j \in \{0, 1\}.
\]
Thus, the $2n$ elements of $D_{2n}$ partition into:
\begin{itemize}\setlength{\itemsep}{1pt}
    \item \textbf{Rotations:} $r^i$ for $0 \le i < n$ (the cyclic subgroup $\langle r \rangle \cong \mathbb{Z}_n$),
    \item \textbf{Reflections:} $r^i s$ for $0 \le i < n$ (elements of order 2).
\end{itemize}

\subsection*{2. Characterization of the Center $Z(D_{2n})$}
The center of a group $G$ is the subgroup of elements commuting with every element in $G$:
\[
Z(D_{2n}) = \big\{ z \in D_{2n} \;\big|\; zg = gz \quad \forall\, g \in D_{2n} \big\}.
\]
Since $D_{2n} = \langle r, s \rangle$, an element $z \in D_{2n}$ belongs to $Z(D_{2n})$ if and only if it commutes with both generators $r$ and $s$:
\[
z \in Z(D_{2n}) \iff zr = rz \quad \text{and} \quad zs = sz.
\]

\subsection*{3. Exhaustive Case Analysis}

\subsubsection*{Case I: Elements of the form $z = r^i s$ (Reflections)}
Let $z = r^i s$ for some $0 \le i < n$. We test whether $z$ commutes with the generator $r$:
\[
zr = (r^i s)r = r^i(sr) = r^i(r^{-1}s) = r^{i-1}s,
\]
\[
rz = r(r^i s) = r^{i+1}s.
\]
For $z \in Z(D_{2n})$, we require $zr = rz$:
\[
r^{i-1}s = r^{i+1}s \implies r^{i-1} = r^{i+1} \implies r^2 = 1 \implies n \mid 2.
\]
However, by hypothesis, $n \ge 3$, which strictly precludes $n \mid 2$. Thus $r^2 \ne 1$, which implies:
\[
zr \ne rz \quad \text{for all } 0 \le i < n.
\]
\textbf{Conclusion for Case I:} No reflection $r^i s$ belongs to the center $Z(D_{2n})$.

\subsubsection*{Case II: Elements of the form $z = r^i$ (Rotations)}
Let $z = r^i$ for some $0 \le i < n$.
\begin{enumerate}
    \item \textbf{Commutation with $r$:}
    \[
    zr = r^i r = r^{i+1} = r r^i = rz.
    \]
    Thus, every rotation $r^i$ commutes with $r$.

    \item \textbf{Commutation with $s$:}
    \[
    zs = r^i s, \qquad sz = s r^i = r^{-i} s \quad \text{(using Eq.~\eqref{eq:rel})}.
    \]
    For $z \in Z(D_{2n})$, we require $zs = sz$:
    \[
    r^i s = r^{-i} s \iff r^i = r^{-i} \iff r^{2i} = 1 \iff 2i \equiv 0 \pmod n \iff n \mid 2i.
    \]
\end{enumerate}

Since $0 \le i < n$, we have $0 \le 2i < 2n$. Therefore, the only possible multiples of $n$ for $2i$ in this range are $2i = 0$ or $2i = n$:

\begin{itemize}
    \item \textbf{If $2i = 0$:}  
    This gives $i = 0$, corresponding to the identity element $z = r^0 = 1 \in Z(D_{2n})$.

    \item \textbf{If $2i = n$:}
    \begin{itemize}
        \item \textbf{Part (a) --- $n$ is odd:}  
        If $n$ is odd, then $\gcd(2, n) = 1$, so $n \mid 2i \implies n \mid i$. Since $0 \le i < n$, the only solution is $i = 0$. Hence, no non-identity rotation commutes with $s$.
        \[
        \mathbf{Z(D_{2n}) = \{1\} = 1 \quad \text{if } n \text{ is odd.}}
        \]

        \item \textbf{Part (b) --- $n$ is even ($n = 2k$ for an integer $k \ge 2$):}  
        Substituting $n = 2k$ into $2i = n$ gives:
        \[
        2i = 2k \implies i = k.
        \]
        Since $1 \le k < 2k = n$, this yields a unique non-identity element $z = r^k$.
        Geometrically, $r^k$ represents a half-turn ($180^\circ$ rotation) which inverts all vectors in the plane and thus commutes with every reflection and rotation.
        \[
        \mathbf{Z(D_{2n}) = \{1, r^k\} \quad \text{if } n = 2k.}
        \]
    \end{itemize}
\end{itemize}

\subsection*{4. Summary of Results}
Combining both cases:
\begin{enumerate}
    \item[\textbf{(a)}] If $n$ is odd, $Z(D_{2n}) = 1$.
    \item[\textbf{(b)}] If $n = 2k$, $Z(D_{2n}) = \{1, r^k\}$.
\end{enumerate}

\begin{flushright}
$\blacksquare$
\end{flushright}

\end{document}
